% Generated draft: requires mathematical and source review.
\documentclass[11pt]{article}
\usepackage[margin=25mm]{geometry}
\usepackage{fontspec}
\setmainfont{Latin Modern Roman}
\usepackage{amsmath,amssymb,mathtools,unicode-math}
\setmathfont{Latin Modern Math}
\usepackage{enumitem,needspace,xurl,hyperref}
\setlength{\parindent}{0pt}
\setlength{\parskip}{5pt}
\setlength{\emergencystretch}{3em}
\newcommand{\N}{\mathbb{N}}
\newcommand{\Z}{\mathbb{Z}}
\newcommand{\Q}{\mathbb{Q}}
\newcommand{\R}{\mathbb{R}}
\newcommand{\C}{\mathbb{C}}
\newcommand{\F}{\mathbb{F}}
\newcommand{\A}{\mathbb{A}}
\newcommand{\PP}{\mathbb{P}}
\newcommand{\E}{\mathbb{E}}
\newcommand{\eps}{\varepsilon}
\newcommand{\dd}{\,\mathrm{d}}
\begin{document}
\section*{UnsolvedMath ID: 6700070}
\textbf{Source problem code:} catalogue entry
\subsection{Definitions and mathematical statement}
Let $(X,g)$ be a complete oriented Riemannian $n$-manifold and let $S^n(1)$ be the unit round sphere. The source defines $\Psi(X)$ as the space of degree-one maps $f:X\to S^n(1)$ that are $1$-Lipschitz and locally constant at infinity. The Lipschitz condition means
\[
d_{S^n}(f(x),f(y))\leq d_X(x,y)\qquad(x,y\in X).
\]
Locally constant at infinity ordinarily means locally constant off a compact subset; on multiple ends this need not mean one common constant outside that subset. The exact end and degree conventions are not reproduced.

The parametric hypersphericity conjecture states that, for $m\geq2$ and $\varepsilon>0$,
\[
\operatorname{Sc}_g\geq m(m-1)+\varepsilon
\quad\Longrightarrow\quad
\operatorname{dim}_{\mathrm{macro}}\Psi(X)\leq n-m-1.
\]
The sphere comparison parameter $m$ is naturally integral, but the source does not give its range relative to $n$. Degree is exactly one, not merely nonzero, and no spin hypothesis is supplied.

Macroscopic dimension is a large-scale metric notion; in one usual convention it is the least dimension of a polyhedron admitting an appropriate continuous uniformly cobounded map from the space. The source gives no metric, topology, coboundedness convention, or definition of macroscopic dimension on $\Psi(X)$. This cannot be replaced automatically by the uniform-distance metric: since the target sphere is bounded, that metric makes every mapping space bounded and loses the intended large-scale information. The fragment's notation $\operatorname{Lip}_\lambda$ also contains an unexplained parameter while its words specify $1$-Lipschitz maps.

For noncompact $X$, defining degree for locally constant-at-infinity maps requires an appropriate compactification or compact-support convention; ordinary proper degree into the compact sphere cannot be assumed. When $n-m-1<0$, an empty-space/negative-dimension convention or a restriction on $m$ is also needed. These missing data are essential to a precise interpretation of the displayed bound.
\subsection{Short English statement}
A scalar-curvature lower bound should force a dimension drop in the large-scale geometry of degree-one short maps to the sphere. The required mapping-space metric, degree-at-infinity and negative-dimension conventions are missing.
\subsection{Sources}
\begin{itemize}
\item Local UnsolvedMath registry, permanent ID 6700070, exact source code catalogue entry, ``Parametric Hypersphericity''; proposed by Misha Gromov, 2017. Registry status open, classification OPEN-TRIAGE. Source printed tag [?74], occurrence 73.
\item Misha Gromov, \emph{101 Questions, Problems and Conjectures around Scalar Curvature} (2017), PDF page 72, \url{https://www.ihes.fr/~gromov/wp-content/uploads/2018/08/101-problemsOct1-2017.pdf}. Supplied normalized pdf-text extraction accessed July 29, 2026; incomplete author-hosted PDF. The marked footnote 66 is not included. No PDF was accessed here.
\item The supplied summary relates the conjecture to hypersphericity and macroscopic-dimension programmes and reports no dedicated resolution. It gives no precise partial-result references or mapping-space definitions.
\item Provenance flags current status and statement-reuse terms NEEDS\_REVIEW; the source marks the passage as a numbered open item.
\end{itemize}
\end{document}
